\section{Linear regression} \label{sec:linear_regression}
Assume the following construction:
\[y_i = \beta^T x_i + \varepsilon_i, \quad \varepsilon_1, ..., \varepsilon_n \overset{iid.}{\sim} \mathcal{N}(0, \sigma^2)\]
where $x_i \in \R^d$ and $\beta \in \R^d$ are fixed

\subsection{Preliminaries}
\begin{lemma}                                                                      
    Let $X: \Omega \to \R^d$ be a random variable and $A \in \R^{d \times d}, b \in \R^d$. Assume that the second moments exist, i.e. $\E [X_i^2] < \infty$ and the covariance matrix is given by $\Sigma$. Then the following properties hold:
    \begin{enumerate}[label=Property \thelemma.\arabic*, align=left]
        \item \label{prop:expectation linearity} $\mathbb{E}[AX + b] = A \cdot \mathbb{E}[X] + b$
        \item \label{prop:covariance linearity} $\cov (AX) = A \cdot \Sigma \cdot A^T$
    \end{enumerate}
\end{lemma}
\begin{remark}                                                    
    Using the fact that the second moment exists and by Jensen's inequality, we have that $\E[X_i] < \infty$ and $\E[X_i^2] < \infty$ for all $i = 1, ..., d$. Therefore the covariance matrix is well-defined, since $\left(\E [|X_i X_j|]\right)^2 = (\E [|X_i| |X_j|])^2 \leq \E [|X_i|^2] \E [|X_j|^2] < \infty$ by Cauchy-Schwarz inequality.
\end{remark}
\begin{proof} \hfill  
    \begin{enumerate}                                                                                                                                        
        \item Trivial
        \item Suppose that the expected value of $X$ is the zero vector. By \ref{prop:expectation linearity} we also have that $\E[AX] = A \E[X] = 0$. Then we have that:
              \begin{align*}                                                                                                                                            
                  \cov((AX)_i, (AX)_j)
                   & = \cov \left(\sum_{k=1}^{d} A_{ik} X_k, \sum_{l=1}^{d} A_{jl} X_l\right)
                  = \E \left(\left(\sum_{k=1}^{d} A_{ik} X_k\right)\left(\sum_{l=1}^{d} A_{jl} X_l\right) \right)
                  \\ & = \E \left[\sum_{k=1}^{d} \sum_{l=1}^{d} A_{ik} X_k X_l A^T_{lj} \right]
                  = \sum_{k=1}^{d} \sum_{l=1}^{d} A_{ik} \E [X_k X_l] A^T_{lj}
                  \\ & = \sum_{k=1}^{d} A_{ik} \left(\sum_{l=1}^{d} \Sigma_{kl} A^T_{lj} \right)
                  = \sum_{k=1}^{d} A_{ik} (\Sigma A^T)_{kj}
                  = (A \Sigma A^T)_{ij}
              \end{align*}
              
              If the expected value of $X$ is not the zero vector, we can define a new random variable $Y = X - \E[X]$, which has the expected value of zero. Then we have that:
              \[\cov (Y_i, Y_j)
                  = \E [Y_i Y_j]
                  = \E[(X_i - \E[X]_i)(X_j - \E[X]_j)]
                  = \cov (X_i, X_j)\]
              Therefore we have that:
              \begin{align*}                                                                                                                                    
                  \cov((AX)_i, (AX)_j)
                   & = \E[(AX - \E[AX])_i \cdot (AX - \E[AX])_j]
                  \\ &\overset{*}{=} \E[(A(X - \E[X]))_i \cdot (A(X - \E[X]))_j]
                  \\ &= \E[(AY)_i \cdot (AY)_j]
                  = \cov((AY)_i, (AY)_j) 
                  \\ &= (A \cov(Y) A^T)_{ij} 
                  = (A \cov(X) A^T)_{ij}
              \end{align*}
              where $*$ follows from \ref{prop:expectation linearity}.
    \end{enumerate}
\end{proof} 

\begin{lemma}[Multivariate normal distribution] \label{lem:multinormal}
    We say that $X: \Omega \to \R^d$ is multivariately normally distributed if and only if there exists a matrix $A \in \R^{d \times d}$ and a vector $b \in \R^d$ such that $X = A Z + b$, where $Z_1, ..., Z_d \overset{iid.}{\sim} \mathcal{N}(0, 1)$.
    
    We have that multivariate normal distribution is closed under linear transformations, i.e. if $X \sim \mathcal{N}(\mu, \Sigma)$ and $Y = AX + b$, then $Y \sim \mathcal{N}(A\mu + b, A \Sigma A^T)$. (For the expectation and covariance refer to \ref{prop:expectation linearity} and \ref{prop:covariance linearity}).
\end{lemma}
\begin{lemma} \label{lem:multivariate independence}
    For the multivariate normal distribution, we have that the components are independent if and only if they are uncorrelated, i.e. $\cov(X_i, X_j) = 0$ for all $i \neq j$.
\end{lemma}

\begin{lemma} \label{lem:projection properties}
    Let $X \in \R^{n \times d}$ be a matrix of full rank, where $n > d$. Define $H = X (X^TX)^{-1}X^T \in \R^{n \times n}$ and $E = \diag \{\underbrace{1, ..., 1}_{n - d}, \underbrace{0, ..., 0}_{d}\}$, then:
    \begin{enumerate}[label=Property \thelemma.\arabic*, align=left]
        \item \label{prop:idempotent symmetric} $I - H$ is symmetric and idempotent.
        \item \label{prop:rank} $\rank(I - H) = n - d$
        \item \label{prop:identity} There exists a matrix $Q$ such that $I - H = Q^T E Q$
        \item \label{prop:kernel image} $\ker (I - H) = \im (X)$
    \end{enumerate}
\end{lemma}
\begin{proof} \hfill 
    \begin{enumerate}                                                                                                                    
        \item This follows by the direct computation:
              \[(I - H)^T 
                  = I - H^T 
                  = I - (X (X^TX)^{-1}X^T)^T
                  = I - (X^T)^T \left((X^TX)^T\right)^{-1} X^T
                  = I - H\]
              \begin{align*}                                                     
                  (I - H)(I - H)
                   & = (I - X (X^TX)^{-1}X^T)(I - X (X^TX)^{-1}X^T)
                  \\ & = I - 2 X (X^TX)^{-1}X^T + X (X^TX)^{-1}X^T X (X^TX)^{-1}X^T
                  \\ & = I - 2 X (X^TX)^{-1}X^T + X (X^TX)^{-1} X^T                 
                  \\ & = I - X (X^TX)^{-1}X^T
                  = I - H
              \end{align*}
        \item Since $I - H$ is idempotent we have that its trace is equal to its rank. Therefore we have that:
              \[\tr(H) 
                  = \tr(X (X^TX)^{-1}X^T)
                  \overset{*}{=} \tr((X^TX)^{-1} X^T X)
                  = \tr(I_d)
                  = d\]
              where $*$ follows from the cyclic property of the trace. Therefore we have that:
              \[\rank(I - H) = \tr(I - H) = \tr(I) - \tr(H) = n - d\]
        \item Since the matrix $I - H$ is idempotent we have that it has an eigenvalue $1$ with multiplicity equal to $\rank(I - H) = n - d$ and eigenvalue $0$ with multiplicity $d$. Its minimal polynomial is contained in $\{x, x-1, x(x-1)\}$ implying that $I - H$ is diagonalizable since in all the cases it does not have repeated factors. Therefore there exists a matrix $Q \in O(\R^n)$, s.t. 
              \[Q^T (I - H) Q = E\]
        \item By the rank-nullity theorem we have $\dim \ker(I - H) = n - \rank(I - H) = n - (n - d) = d$. Using the assumption we have that $\dim \im(X) = d$. Therefore it suffices to show that $\im X \subset \ker(I - H)$:
              \begin{align*}                                
                  (I - H)Xy 
                   & = (I - X(X^TX)^{-1}X^T)Xy 
                  \\ &= Xy - X(X^TX)^{-1}X^TXy
                  = Xy - Xy = 0 \quad \forall y \in \R^d
              \end{align*}
              This shows that $\im X \subset \ker(I - H)$, and since they have the same dimension we have that $\im X = \ker(I - H)$.
    \end{enumerate}
\end{proof}

\subsection{t-statitistic}
In this section we assume that $\hat{\beta}$ is the ordinary least squares estimator of $\beta$. We establish two facts:
\begin{itemize}                                        
    \item $\varepsilon \sim \mathcal{N}(0, \sigma^2 I)$ since the components of $\varepsilon$ are independent and identically distributed as $\mathcal{N}(0, \sigma^2)$, thus we can obtain $\varepsilon = \sigma I \cdot (\varepsilon / \sigma)$, where $\varepsilon_1 / \sigma, ..., \varepsilon_n / \sigma \overset{iid.}{\sim} \mathcal{N}(0, 1)$.
    \item $y = X\beta + \varepsilon \sim \mathcal{N}(X\beta, \sigma^2 I)$ as a consequence of Lemma~\ref{lem:multinormal} and the fact that $y$ is a linear transformation of $\varepsilon$.
\end{itemize}
\begin{enumerate}[label=Step \arabic*., align=left]
    \item Compute the distribution of the residual sum of squares: 
          \[\hat{\varepsilon}
              = y - X \hat{\beta}
              = y - X (X^TX)^{-1}X^T y
              = (I - H) y\]
          
          By \ref{prop:kernel image} we have that $\im(X) = \ker(I - H)$, and since $X\beta \in \im(X)$ we have that:
          \[(I - H) X\beta = 0\]
          thus:
          \[(I - H) y 
              = (I - H) (X\beta + \varepsilon) 
              = (I - H)\varepsilon\]
          with the notation from Lemma~\ref{lem:projection properties}. Using the result from \ref{prop:identity} we have that:
          \[I - H = Q^T E Q\]
          for some orthogonal matrix $Q$.
          
          Writing out the expression for the residual sum of squares, we have that:
          \[\text{RSS} 
              = \hat{\varepsilon}^T \hat{\varepsilon}
              = \varepsilon^T (I - H)^T (I - H) \varepsilon
              = \varepsilon^T (I - H) \varepsilon\]
          
          Define $U = Q \varepsilon$. Then $U \sim \mathcal{N}(0, \sigma^2 I)$ by Lemma~\ref{lem:multinormal}. Then we have that:
          \begin{itemize}                                                                                    
              \item $\E[Q \varepsilon] = Q \E[\varepsilon] = 0$ by \ref{prop:expectation linearity}
              \item $\cov(Q \varepsilon) = Q \cov(\varepsilon) Q^T = Q \sigma^2 I Q^T = \sigma^2 Q Q^T = \sigma^2 I$ by \ref{prop:covariance linearity}
          \end{itemize}
          
          Then by Lemma~\ref{lem:multivariate independence} we have that the components of $U$ are independent, since they are uncorrelated.
          
          Thus we have that $U/\sigma$ follows a joint standard normal distribution. Then we have that:
          \begin{equation} \label{eq:rss expansion}
              \text{RSS} = U^T E U
          \end{equation}
          \[\frac{\text{RSS}}{\sigma^2}
              = U^T / \sigma \cdot E \cdot U / \sigma
              = \sum_{i=1}^{n - d} (U_i / \sigma)^2
              \sim \chi^2_{n - d}\]
    \item Compute the distribution of $\hat{\beta}_i - \beta_i$:
          
          Observe that the $\hat{\beta}$ follows multivariate normal distribution since $\hat{\beta} = Ay$ where $A = (X^TX)^{-1}X^T$, and $y \sim \mathcal{N}(X\beta, \sigma^2 I)$ (see Lemma~\ref{lem:multinormal}). 
          Therefore $\hat{\beta} \sim \mathcal{N} (\beta, \sigma^2 AA^T)$, where we have used that $\hat{\beta}$ is an unbiased estimator of $\beta$ and that
          \[\cov(\hat{\beta}) 
              = \cov(Ay)
              = A\cov(y)A^T
              = \sigma^2 AA^T\]
          by \ref{prop:covariance linearity}.
          \[AA^T 
                  = (X^TX)^{-1} X^T X \left((X^TX)^{-1}\right)^T
              = \left((X^TX)^{-1}\right)^T
              = \left((X^TX)^T\right)^{-1}
              = (X^TX)^{-1}\]
          
          As was derived above, we have that $\hat{\beta} \sim \mathcal{N} (\beta, \sigma^2 (X^TX)^{-1})$. Therefore we have that:
          \[\hat{\beta}_i - \beta_i \sim \mathcal{N}(0, \sigma^2 (X^TX)^{-1}_{ii})\]
          
          Thus:
          \[\frac{\hat{\beta}_i - \beta_i}{\sigma \sqrt{(X^TX)^{-1}_{ii}}} \sim \mathcal{N}(0, 1)\]
    \item $\hat{\varepsilon}$ and $\hat{\beta} - \beta$ are independent:
          
          We write $\hat{\varepsilon} = (I - H) \varepsilon$ and $\hat{\beta} - \beta = (X^TX)^{-1} X^T (X\beta + \varepsilon) - \beta = (X^TX)^{-1} X^T \varepsilon$. Then we have that:
          
          \[\begin{pmatrix}                    
                  \hat{\varepsilon} \\ 
                  \hat{\beta} - \beta
              \end{pmatrix} =
              \begin{pmatrix}                    
                  I - H \\
                  (X^TX)^{-1} X^T
              \end{pmatrix} \varepsilon\]
          
          Thus by Lemma~\ref{lem:multinormal} we have that the joint distribution of $\hat{\varepsilon}$ and $\hat{\beta} - \beta$ is multivariate normal. Therefore by Lemma~\ref{lem:multivariate independence} it suffices to show that they are uncorrelated, i.e. $\cov(\hat{\varepsilon}, \hat{\beta} - \beta) = 0$. We have that:          
          \begin{align*}                
              \cov (\hat{\varepsilon}, \hat{\beta} - \beta)
              \overset{*}{=} \E [\hat{\varepsilon} (\hat{\beta} - \beta)^T]
               & = \E [(I - H) \varepsilon \varepsilon^T X (X^TX)^{-1}]
              \\ &\overset{**}{=} (I - H) \E [\varepsilon \varepsilon^T] X (X^TX)^{-1}
              \\ &= \sigma^2 (I - H) X (X^TX)^{-1}
          \end{align*}
          where $*$ follows from the fact that both $\hat{\varepsilon}$ and $\hat{\beta} - \beta$ have expected value of zero as the linear transformation of $\varepsilon$ which has expected value of zero. $**$ follows from \ref{prop:expectation linearity}.
          
          By Lemma~\ref{prop:kernel image} we have that $\im(X) = \ker(I - H)$, thus $(I - H) X = 0$, hence:
          \[\cov (\hat{\varepsilon}, \hat{\beta} - \beta)
              = \sigma^2 (I - H) X (X^TX)^{-1}
              = 0\]
    \item Conclude that the t-statistic follows a t-distribution with $n - d$ degrees of freedom:
          \[\frac{\hat{\beta}_i - \beta_i}{\sqrt{\frac{\text{RSS}}{n - d} (X^TX)^{-1}_{ii}}}
              = \frac{\frac{\hat{\beta}_i - \beta_i}{\sigma \sqrt{(X^TX)^{-1}_{ii}}}}{\frac{1}{\sigma} \sqrt{\frac{\text{RSS}}{n - d}}}
              = \frac{\frac{\hat{\beta}_i - \beta_i}{\sigma \sqrt{(X^TX)^{-1}_{ii}}}}{\sqrt{\frac{\text{RSS}}{\sigma^2 (n - d)}}}
              = \frac{Z}{\sqrt{\frac{T}{n - d}}}\]
          where by the previous steps we have that $Z = \frac{\hat{\beta}_i - \beta_i}{\sigma \sqrt{(X^TX)^{-1}_{ii}}} \sim \mathcal{N}(0, 1)$ and $T = \frac{\text{RSS}}{\sigma^2} \sim \chi^2_{n - d}$. $T$ depends only on the residuals and $Z$ depends only on $\hat{\beta}_i$, which are independent by the previous step. Therefore we have that the t-statistic follows a t-distribution with $n - d$ degrees of freedom.
\end{enumerate}

\begin{theorem} \label{thm:variance estimator}
    The unbiased estimator of the variance of the error term is given by:
    \[\hat{\sigma}^2 = \frac{1}{n - d} \sum_{i=1}^{n} \hat{\varepsilon}_i^2\]
    where $d$ is the number of features and $n$ is the number of samples.    
\end{theorem}
\begin{proof}                                                             
    \[\text{RSS} 
        = \hat{\varepsilon}^T\hat{\varepsilon}
        \overset{\eqref{eq:rss expansion}}{=} \sum_{i=1}^{n - d} U_i^2\]
    
    Each $U_i$ is distributed as $\mathcal{N}(0, \sigma^2)$, therefore we have that:
    \[\E[U_i^2] = \var(U_i) = \sigma^2\]
    
    Thus we have that:
    \[\E[\text{RSS}] = \E\left[\sum_{i=1}^{n - d} U_i^2\right] = \sum_{i=1}^{n - d} \E[U_i^2] = (n - d) \sigma^2\]
    
    Dividing both sides by $n - d$ yields:
    \[\E\left[\frac{\text{RSS}}{n - d}\right] = \sigma^2\]
\end{proof}
\begin{remark}                                                
    This provides an intuition why we use this estimator for the variance of the error term.
\end{remark}

\begin{theorem}[Gauss-Markov theorem] \label{thm:gauss-markov}
    The ordinary least squares estimator $\hat{\beta}$ is the best linear unbiased estimator (BLUE) of $\beta$, i.e. for each $C \in \R^{d \times n}$ such that $\E[C y] = \beta$, we have that:
    \[\cov(\hat{\beta}) \preceq \cov(Cy)\]    
\end{theorem}
\begin{proof}            
    We need to show that the matrix:
    \[\cov(Cy) - \cov(\hat{\beta})\]
    is positive semidefinite.
    
    Define $D = C - (X^TX)^{-1}X^T$. Then we have that:
    \[\E [Cy] 
        = \E [((X^TX)^{-1}X^T + D)y]
        = \E [(X^TX)^{-1}X^Ty] + \E[Dy]
        = \beta\]
    
    Thus we have that $\E[Dy] = 0$. Therefore we have that:
    \[\E[Dy] 
        = \E[DX\beta + D\varepsilon]
        = DX\beta = 0\]
    
    Since $\beta$ is unknown we must have that $DX\beta = 0$ for all $\beta$, which implies that $DX = 0$. This also implies $(DX)^T = X^TD^T = 0$. Then we have that:
    \begin{align*}    
        \cov(Cy)
        &= \sigma^2 CC^T
        \\ &= \sigma^2 ((X^TX)^{-1}X^T + D)((X^TX)^{-1}X^T + D)^T
        \\ &= \sigma^2 ((X^TX)^{-1}X^TX(X^TX)^{-1} + DX(X^TX)^{-1} + (X^TX)^{-1}X^TD^T + DD^T)
        \\ &= \sigma^2 ((X^TX)^{-1} + DD^T)
    \end{align*}

    We know that $DD^T$ is positive semidefinite, so for any $v \in \R^d$ we have that:
    \[v^T CC^T v = v^T (X^TX)^{-1} v + v^T DD^T v \geq v^T (X^TX)^{-1} v\]
    thus:
    \[v^T (CC^T - (X^TX)^{-1}) v \geq 0\]
    implying that $CC^T - (X^TX)^{-1}$ is positive semidefinite. Therefore we have that:
    \[\cov(Cy) - \cov(\hat{\beta}) = \sigma^2 (CC^T - (X^TX)^{-1}) \succeq 0\]
\end{proof}